\section{Literature Review}\label{sec:litreview}

\subsection{SGP4 Accuracy Collapse During Storm-Time Drag}
Storm-time drag modeling remains the primary failure mode for catalog-grade LEO propagation. Thermospheric density enhancements under strong geomagnetic forcing produce rapid semi-major-axis decay and along-track phase drift that are not captured by static $B^*$ assumptions \cite{oliveira2021current,walter2025orbital}. Event studies centered on the May 2024 storm confirm that unmodeled density surges can produce tens-of-kilometers errors and temporary tracking degradation in dense constellations \cite{parker2024satellite,kang2025solar}. Data-assimilative drag updates improve this behavior, but are not yet universally integrated in operational SGP4 pipelines \cite{dey2024improving}.

\subsection{PINNs, Orbit Determination, and Uncertainty (2025--2026)}
Recent work has moved PINNs from proof-of-concept toward practical astrodynamics use. PINN-based orbit determination studies report robust fitting under sparse observations and physically plausible trajectories \cite{loshelder2025od}. In parallel, covariance-aware PINN formulations provide uncertainty outputs suitable for conjunction workflows \cite{dallinger2025cov}. Complementary ML thermospheric models now expose dynamic confidence bounds, enabling risk assessments that reflect environment uncertainty rather than deterministic single-track forecasts \cite{li2025density,mutschler2025density}.

\subsection{Hybrid Models for Mega-Constellation Operations}
Hybrid residual-correction architectures remain attractive for scale because they preserve SGP4 compatibility while learning regime-specific biases. For OneWeb-like shells, ML-assisted methods substantially outperform open-loop SGP4 for multi-day horizons, but reliability under extreme storms is still a limiting factor \cite{chen2025oneweb,bertolini2024hybrid}. Broader space-safety studies increasingly advocate mixed analytical--ML stacks with explicit uncertainty handling and verification constraints for deployment in operational mission planning \cite{rommel2025verifiable,ramesh2025space}. Our G-PINN follows this trajectory by combining hard kinematic constraints, weather-aware features, and probabilistic conjunction metrics in one deployable framework.
